\subsection{Operador Ternario}

\subsubsection{Regla: solo para seleccionar un valor simple}

\textbf{Regla:} El operador condicional \texttt{?:} se usa cuando una condición sencilla selecciona entre dos valores sin efectos secundarios. Se prohíben operadores ternarios anidados; cuando existen más decisiones o acciones, se utiliza una estructura \texttt{if}/\texttt{else} (Microsoft, 2026e).

\textbf{Código correcto:}
\begin{lstlisting}[style=csharpstyle]
string status = player.IsActive ? "Active" : "Inactive";
\end{lstlisting}

\textbf{Código incorrecto:}
\begin{lstlisting}[style=csharpstyle]
string status = player.IsActive ? player.HasWon ? "Winner" : "Playing" : "Inactive";
\end{lstlisting}
